% !PRINTED 422
\weberpage{422}

%<<CONT
Co\"efficienten von $x_1^{11}x_2,\ x_1^{22}x_2^2,\ x_1^{30}$ auf beiden Seiten der Gleichungen (14), und findet so die Relation
\weq{15}{t^5 - 10ft^3 + 45f^2t - T = 0,}
oder auch
\weq{16}{t^2(t^4 - 10ft^2 + 45f^2)^2 - T^2 = 0.}

Aus (16) und (15) ergeben sich nun nach (10) und (12) die Gleichungen 5\textsuperscript{ten} Grades f\"ur $r$ und f\"ur $u$
\weq{17}{r(r^2 - 10r + 45)^2 = z_1,}
\weq{18}{u^5 - \frac{10u^3}{z_1} + \frac{45u}{z_1^2} - \frac{1}{z_1^2} = 0.}

Die letzte dieser Gleichungen geht durch die Substitution
\weq{19}{y = -\frac{1}{3u},\quad z_1 = 27\gamma}
in die Form
\weq{20}{y^5 + 15y^4 - 10\gamma y^2 + 3\gamma^2 = 0}
\"uber, die wir bereits im \S.\,75 des ersten Bandes als eine Normalform der Gleichung 5\textsuperscript{ten} Grades kennen gelernt haben.

\begin{center}
\S.\,112.\\[4pt]
\gesperrt{Die Hauptresolvente f\"unften Grades.}
\end{center}

Wenn man die Functionen $u, v$ gleichzeitig benutzt, so kann man eine Schaar von Resolventen ableiten, die die Form der Hauptgleichung 5\textsuperscript{ten} Grades haben, und sich direct mit einer beliebig gegebenen Hauptgleichung 5\textsuperscript{ten} Grades in Uebereinstimmung bringen lassen\footnote{\gesperrt{Kiepert}, G\"ottinger Nachrichten 1878. Crelle's Journal, Bd.\,87. \gesperrt{Klein}, Ikosa\"eder, S.\,106.}.

Da es beim Ikosa\"eder keine Invarianten 8\textsuperscript{ten}, 14\textsuperscript{ten}, 16\textsuperscript{ten}, 22\textsuperscript{sten} oder 28\textsuperscript{sten} Grades giebt, so m\"ussen die Functionen $S(\tilde\omega)$, $S(t\tilde\omega)$, $S(\tilde\omega^2)$, $S(t\tilde\omega^2)$, $S(t^2\tilde\omega^2)$ verschwinden, und wenn wir also
\weq{1}{Y = \alpha\tilde\omega + \beta t\tilde\omega}
setzen, so werden auch die Functionen
\[ S(Y),\quad S(Y^2) \]
f\"ur beliebige Werthe der Parameter $\alpha, \beta$ verschwinden, und daraus ergiebt sich eine identische Gleichung von der Form
%CONT>>
